\section{Appendix}

\subsection{Website Implementation}\label{app:site_implementation}
Given the selected websites described in \S\ref{sec:website_selection}, we make the best attempt to reproduce the functionality of commonly used sites in a reproducible way. 
To achieve this, we utilized open-source frameworks for the development of the websites across various categories and imported data from their real-world counterparts.
For the E-commerce category, we constructed a shopping website with approximately 90$k$ products, including the prices, options, detailed product descriptions, images, and reviews, spanning over 300 product categories. This website is developed using Adobe Magento, an open-source e-commerce platform\footnote{\url{https://github.com/magento/magento2}}. Data resources were obtained from data from actual online sites, such as that included in the Webshop data dump\cite{yao2022webshop}.
As for the social forum platform, we deployed an open-source software Postmill\footnote{\url{https://postmill.xyz/}}, the open-sourced counterpart of Reddit\footnote{\url{https://www.reddit.com/}}.
We sampled from the top 50 subreddits\footnote{\url{https://redditlist.com/sfw.html}}. 
We then manually selected many subreddit for northeast US cities as well as subreddit for machine learning and deep learning-related topics. This manual selection encourages cross-website tasks such as seeking information related to the northeast US on both Reddit and the map. 
In total, we have 95 subreddits, 127390 posts, and 661781 users. 
For the collaborative software development platform, we choose GitLab\footnote{\url{https://gitlab.com/gitlab-org/gitlab}}.
%\dfried{did you reimplement all of GitLab? or just some important parts} \gn{It's not really reimplement, right, you use the open-source.} 
We heuristically simulate the code repository characteristics by sampling at least ten repositories for every programming language: $80\%$ of them are sampled from the set of top $90$ percentile wrt stars repos using a discrete probability distribution weighted proportional to their number of stars; the remaining are sampled from the bottom ten percentile set using similar weighted distribution. This is done to ensure fair representation of repos of all kinds, from popular projects with many issues and pull requests to small personal projects. In total, we have 300 repositories and more than 1000 accounts with at least one commit to a repository.
For the content management system, we adapted Adobe Magento's admin portal, deploying the sample data provided in the official guide.
We employ OpenStreetMap\footnote{\url{https://www.openstreetmap.org/}} for map service implementation, confining our focus to the northeast US region due to data storage constraints.
We implement a calculator and a scratchpad ourselves.

Lastly, we configure the knowledge resources as individual websites, complemented with search functionality for efficient information retrieval.
Specifically, we utilize Kiwix\footnote{\url{https://www.kiwix.org/en/}} to host an offline version of English Wikipedia with a knowledge cutoff of May 2023. 
The user manuals for GitLab and Adobe Commerce Merchant documentation are scraped from the official websites.%\sz{Not too interesting and will break the reading flow, maybe summarize and move to appendix} 

\subsection{Environment Delivery and Reset}\label{app:env_reset}
One goal for our evaluation environment is ease of use and reproducibility.
As a result, we deploy our websites in separate Docker images~\footnote{\url{https://www.docker.com/}}, one per website.
The Docker images are fully self-contained with all the code of the website, database, as well as any other software dependencies.
They also do not rely on external volume mounts to function, as the data of the websites are also part of the docker image.
This way, the image is easy to distribution containing all the pre-populated websites for reproducible evaluation.
End users can download our packaged Docker images and run them on their systems and re-deploy the exact websites together with the data used in our benchmarks for their local benchmarking.

Since some evaluation cases may require the agent to modify the data contained in the website, \eg creating a new user, deleting a post, etc., it is crucial to be able to easily reset the website environment to its initial state.
With Docker images, the users could stop and delete the currently running containers for that website and start the container from our original image again to fully reset the environment to the initial state.
Depending on the website, this process may take from a few seconds to one minute.
However, not all evaluation cases would require an environment reset, as many of the intents are information gathering and are read-only for the website data.
Also, combined with the inference time cost for the agent LLMs, we argue that this environment reset method, through restarting Docker containers from the original images, will have a non-negligible but small impact on evaluation time.

% \subsection{Annotation Instruction}\label{app:annotation_instruction}
% The full instruction for annotating the intents can be found in \url{}.
\subsection{User Roles Simulation}\label{sec:user_role}
Users of the same website often have disparate experiences due to their distinct \emph{roles}, \emph{permissions}, and \emph{interaction histories}. 
For instance, within an E-commerce CMS, a shop owner might possess full read and write permissions across all content, whereas an employee might only be granted write permissions for products but not for customer data. We aim to emulate this scenario by generating unique user profiles on each platform.

On the shopping site, we created a customer profile that has over 35 orders within a span of two years. 
%\dfried{I think this is the first mention of OneStopShop?}
On GitLab, we selected a user who maintains several popular open-source projects with numerous merge requests and issues. This user also manages a handful of personal projects privately. 
On Reddit, our chosen profile was a user who actively participates in discussions, with many posts and comments. 
Lastly, on our E-commerce CMS, we set up a user profile for a shop owner who has full read-and-write access to all system contents.%\sz{To appendix}
%\dfried{it's a bit confusing since both (1) and (4) are shop-related}

All users are automatically logged into their accounts using a pre-cached cookie. To our best knowledge, this is the first publicly available agent evaluation environment to implement such a characteristic. Existing literature typically operates under the assumption of universally identical user roles~\cite{shi2017world,liu2018reinforcement,deng2023mind2web}. %\dfried{maybe ``user roles'' instead, since intents vary from instruction-to-instruction in these past works}
\subsection{Intent Distribution}
The distribution of intents across the websites are shown in \autoref{fig:intent_distribution}.
\begin{figure}
    \centering
    \includegraphics[width=0.5\linewidth]{figs/site_distribution.pdf}
    \caption{The intent distribution across different websites. Cross-site intents necessitate interacting with multiple websites. Notably, regardless of the website, all user intents require interactions with multiple web pages.}
    \label{fig:intent_distribution}
\end{figure}

\subsection{Human Performance}\label{app:human_annotation}
We acknowledge that there may be a difference in human performance when annotators with different demographics are involved. In fact, many tasks in our dataset require domain-specific knowledge. For instance, an average user may not know what a git merge request is; or how to create a product in a complex content management system. We aim to design tasks that have easy-to-imagine outcomes (\eg a new product page is created) rather than those that are easily performed by an average user without significant domain knowledge.

\subsection{Experiment Configurations}\label{sec:exp_configs}
We experiment with \textsc{GPT-3.5-turbo-16k-0613}, \textsc{GPT-4-0613}, and \textsc{text-bison-001} with a temperature of $1.0$ and a top-$p$ parameter of $0.9$. The maximum number of state transitions is set to 30.
% Notably, we use a high temperature to encourage exploration.
%\dfried{max tokens, or state transitions, or something else?} 
We halt execution if the same action is repeated more than three times on the same observation or if the agent generates three consecutive invalid actions. These situations typically indicate a high likelihood of execution failure and hence warrant early termination. For \textsc{text-bison-001}, we additionally allow ten retries until it generates a valid action.

Primarily, we use a high temperature of 1.0 to encourage the \emph{exploration}. To aid replicating the results, we provide the results of \textsc{GPT-3.5-turbo-16k-0613} with temperature 0.0 in \autoref{tab:temp_zero_results} and the execution trajectories in our code repository.

\begin{table}
    \centering
    \begin{tabular}{@{}c@{\hspace{3pt}}c@{}l@{\hspace{5pt}}c@{}}
        \toprule
        \textbf{CoT} & \textbf{UA Hint} & \multicolumn{1}{c}{\textbf{Model}} & \textbf{SR}  \\
        \midrule
        \cmark & \xmark & \scriptsize\textsc{GPT-3.5} &  6.28 \\
        \bottomrule
    \end{tabular}
    \caption{The task success rate (SR \%) of \textsc{GPT-3.5-turbo-16k-0613} with temperature 0.0.}
    \label{tab:temp_zero_results}
\end{table}


\subsection{Prompt for \texttt{fuzzy\_match}}\label{app:prompt_fuzzy_match}
\fbox{
\parbox{\textwidth}{
Help a teacher to grade the answer of a student given a question. Keep in mind that the student may use different phrasing or wording to answer the question. The goal is to evaluate whether the answer is semantically equivalent to the reference answer.

question: \{\{intent\}\}

reference answer: \{\{reference answer\}\}

all the string 'N/A' that you see is a special sequence that means 'not achievable'

student answer: \{\{prediction\}\}

Conclude the judgement by correct/incorrect/partially correct.
}}

Predictions that are judged as ``correct'' will receive a score of one, while all other predictions will receive a score of zero.

\subsection{The Accuracy of Fuzzy Match Function}\label{app:fuzzy_match_perf}
To evaluate this, we manually checked 40 examples and found that 39 of them are identical to our human judgment. 
In addition, among the 82 examples that require using GPT-4 for evaluation, the answer of 49 (60\%) examples is a date (\eg 10/23/2022) or time duration (\eg 15 minutes). 
In these cases, GPT-4 is only used to judge the different \emph{format} of the answers.
We quantitatively evaluate the correctness of GPT-4 in this case by generating different formats of a date and time duration programmatically. 
We randomly sample negative examples. For instance, Nov 3, 2022, November 3, 2022, 3rd November 2022, 3 Nov 2022, 2022-11-03, and 3rd of November, 2022 are all correct variances of 2022/11/03. The accuracy of GPT-4 is shown in \autoref{tab:fuzzy_match_perf}.
We can see that two versions of GPT-4 are extremely accurate, both achieving 100\% accuracy.
\begin{table}[]
    \centering
    \begin{tabular}{ccc}
    \toprule
     \textbf{Dataset} & \texttt{gpt-4-0613} & \texttt{gpt-4-1106-preview } \\
     \midrule
     Date (900 examples) & 100 & 100 \\
     Time duration (900 examples) & 100 & 100 \\
     \bottomrule
    \end{tabular}
    \caption{The accuracy (\%) of two versions of GPT-4 on judging if dates and time duration of different formats are equivalent.}
    \label{tab:fuzzy_match_perf}
\end{table}


\subsection{The Prompts of the Baseline Web Agents}\label{app:agent_prompts}
The system message of the reasoning agent for both GPT-3.5 and \textsc{GPT-4} is in \autoref{fig:reasoning_prompt}, and two examples are in \autoref{fig:reasoning_examples}. The system message of the direct agent for GPT-3.5 is in \autoref{fig:direct_prompt} and the two examples are in \autoref{fig:direct_example}.
\textbf{UA hint} refers to the instruction of `` If you believe the task is impossible to complete, provide the answer as "N/A" in the bracket.''. We remove this sentence in our ablation studies.

\begin{figure}
\noindent\fbox{\parbox{\textwidth}{%
You are an autonomous intelligent agent tasked with navigating a web browser. You will be given web-based tasks. These tasks will be accomplished through the use of specific actions you can issue.
\\~\\
Here's the information you'll have:

The user's objective: This is the task you're trying to complete.

The current web page's accessibility tree: This is a simplified representation of the webpage, providing key information.

The current web page's URL: This is the page you're currently navigating.

The open tabs: These are the tabs you have open.

The previous action: This is the action you just performed. It may be helpful to track your progress.
\\~\\
The actions you can perform fall into several categories:

Page Operation Actions

\texttt{\`}click [id]\texttt{\`}: This action clicks on an element with a specific id on the webpage.

\texttt{\`}type [id] [content] [press\_enter\_after=0|1]\texttt{\`}: Use this to type the content into the field with id. By default, the "Enter" key is pressed after typing unless press\_enter\_after is set to 0.

\texttt{\`}hover [id]\texttt{\`}: Hover over an element with id.

\texttt{\`}press [key\_comb]\texttt{\`}:  Simulates the pressing of a key combination on the keyboard (e.g., Ctrl+v).

\texttt{\`}scroll [direction=down|up]\texttt{\`}: Scroll the page up or down.
\\~\\
Tab Management Actions:

\texttt{\`}new\_tab\texttt{\`}: Open a new, empty browser tab.

\texttt{\`}tab\_focus [tab\_index]\texttt{\`}: Switch the browser's focus to a specific tab using its index.

\texttt{\`}close\_tab\texttt{\`}: Close the currently active tab.
\\~\\
URL Navigation Actions:

\texttt{\`}goto [url]\texttt{\`}: Navigate to a specific URL.

\texttt{\`}go\_back\texttt{\`}: Navigate to the previously viewed page.

\texttt{\`}go\_forward\texttt{\`}: Navigate to the next page (if a previous 

\texttt{\`}go\_back\texttt{\`} action was performed).
\\~\\
Completion Action:

\texttt{\`}stop [answer]\texttt{\`}: Issue this action when you believe the task is complete. If the objective is to find a text-based answer, provide the answer in the bracket. If you believe the task is impossible to complete, provide the answer as "N/A" in the bracket.
\\~\\
Homepage:

If you want to visit other websites, check out the homepage at http://homepage.com. It has a list of websites you can visit.

http://homepage.com/password.html lists all the account names and passwords for the websites. You can use them to log in to the websites.
\\~\\
To be successful, it is very important to follow the following rules:

1. You should only issue an action that is valid given the current observation

2. You should only issue one action at a time.

3. You should follow the examples to reason step by step and then issue the next action.

4. Generate the action in the correct format. Start with a "In summary, the next action I will perform is`` phrase, followed by action inside \texttt{\`{}\`{}\`{}}\texttt{\`{}\`{}\`{}}. For example, "In summary, the next action I will perform is \texttt{\`{}\`{}\`{}}click [1234]\texttt{\`{}\`{}\`{}}''.

5. Issue stop action when you think you have achieved the objective. Don't generate anything after stop."""
        }
}%
\caption{The \texttt{system} message of the reasoning agent. This message has a general explanation of the task, the available actions, and some notes on avoiding common failures.}\label{fig:reasoning_prompt}
\end{figure}

\begin{figure}
    \centering
    \noindent\fbox{%
    \parbox{\textwidth}{%
\textcolor{red}{example\_user}

OBSERVATION:

[1744] link 'HP CB782A\#ABA 640 Inkjet Fax Machine (Renewed)'

\hspace{2mm}[1749] StaticText '\$279.49'
    
\hspace{2mm}[1757] button 'Add to Cart'
    
\hspace{2mm}[1760] button 'Add to Wish List'
    
\hspace{2mm}[1761] button 'Add to Compare'

URL: http://onestopmarket.com/office-products/office-electronics.html

OBJECTIVE: What is the price of HP Inkjet Fax Machine

PREVIOUS ACTION: None
\\~\\

\textcolor{red}{example\_assistant} 

Let's think step-by-step. This page lists the information of HP Inkjet Fax Machine, which is the product identified in the objective. Its price is \$279.49. I think I have achieved the objective. I will issue the stop action with the answer. In summary, the next action I will perform is \texttt{\`{}\`{}\`{}}stop [\$279.49]\texttt{\`{}\`{}\`{}}
\\~\\

\textcolor{red}{example\_user}

OBSERVATION:

[164] textbox 'Search' focused: True required: False

[171] button 'Go'

[174] link 'Find directions between two points'

[212] heading 'Search Results'

[216] button 'Close'

URL: http://openstreetmap.org

OBJECTIVE: Show me the restaurants near ABC

PREVIOUS ACTION: None
\\~\\
\textcolor{red}{example\_assistant}

Let's think step-by-step. This page has a search box whose ID is [164]. According to the nominatim rule of openstreetmap, I can search for the restaurants near a location by \"restaurants near\". I can submit my typing by pressing the Enter afterwards. In summary, the next action I will perform is \texttt{\`{}\`{}\`{}}type [164] [restaurants near ABC] [1]\texttt{\`{}\`{}\`{}}
        }
}%
\caption{The two examples provided as \texttt{example\_user} and \texttt{example\_assistant} for the reasoning agent. Before issuing the action, the agent first perform reasoning.}
    \label{fig:reasoning_examples}
\end{figure}



\begin{figure}
\noindent\fbox{%
    \parbox{\textwidth}{%
You are an autonomous intelligent agent tasked with navigating a web browser. You will be given web-based tasks. These tasks will be accomplished through the use of specific actions you can issue.
\\~\\
Here's the information you'll have:

The user's objective: This is the task you're trying to complete.

The current web page's accessibility tree: This is a simplified representation of the webpage, providing key information.

The current web page's URL: This is the page you're currently navigating.

The open tabs: These are the tabs you have open.

The previous action: This is the action you just performed. It may be helpful to track your progress.
\\~\\
The actions you can perform fall into several categories:

Page Operation Actions

\texttt{\`}click [id]\texttt{\`}: This action clicks on an element with a specific id on the webpage.

\texttt{\`}type [id] [content] [press\_enter\_after=0|1]\texttt{\`}: Use this to type the content into the field with id. By default, the "Enter" key is pressed after typing unless press\_enter\_after is set to 0.

\texttt{\`}hover [id]\texttt{\`}: Hover over an element with id.

\texttt{\`}press [key\_comb]\texttt{\`}:  Simulates the pressing of a key combination on the keyboard (e.g., Ctrl+v).

\texttt{\`}scroll [direction=down|up]\texttt{\`}: Scroll the page up or down.
\\~\\
Tab Management Actions:

\texttt{\`}new\_tab\texttt{\`}: Open a new, empty browser tab.

\texttt{\`}tab\_focus [tab\_index]\texttt{\`}: Switch the browser's focus to a specific tab using its index.

\texttt{\`}close\_tab\texttt{\`}: Close the currently active tab.
\\~\\
URL Navigation Actions:

\texttt{\`}goto [url]\texttt{\`}: Navigate to a specific URL.

\texttt{\`}go\_back\texttt{\`}: Navigate to the previously viewed page.

\texttt{\`}go\_forward\texttt{\`}: Navigate to the next page (if a previous 

\texttt{\`}go\_back\texttt{\`} action was performed).
\\~\\
Completion Action:

\texttt{\`}stop [answer]\texttt{\`}: Issue this action when you believe the task is complete. If the objective is to find a text-based answer, provide the answer in the bracket. If you believe the task is impossible to complete, provide the answer as "N/A" in the bracket.
\\~\\
Homepage:

If you want to visit other websites, check out the homepage at http://homepage.com. It has a list of websites you can visit.

http://homepage.com/password.html lists all the account name and password for the websites. You can use them to log in to the websites.
\\~\\
To be successful, it is very important to follow the following rules:

To be successful, it is very important to follow the following rules:

1. You should only issue an action that is valid given the current observation

2. You should only issue one action at a time.

3. Generate the action in the correct format. Always put the action inside a pair of \texttt{\`{}\`{}\`{}}. For example, \texttt{\`{}\`{}\`{}}click [1234]\texttt{\`{}\`{}\`{}}

4. Issue stop action when you think you have achieved the objective. Don't generate anything after stop."""
    }%
}
\caption{The \texttt{system} message of the direct agent. This message has the general explanation of the task, the available actions and some notes on avoiding common failures.}\label{fig:direct_prompt}
\end{figure}

\begin{figure}
    \centering
    \noindent\fbox{%
    \parbox{\textwidth}{%
\textcolor{red}{example\_user}

OBSERVATION:

[1744] link 'HP CB782A\#ABA 640 Inkjet Fax Machine (Renewed)'

\hspace{2mm}[1749] StaticText '\$279.49'
    
\hspace{2mm}[1757] button 'Add to Cart'
    
\hspace{2mm}[1760] button 'Add to Wish List'
    
\hspace{2mm}[1761] button 'Add to Compare'

URL: http://onestopmarket.com/office-products/office-electronics.html

OBJECTIVE: What is the price of HP Inkjet Fax Machine

PREVIOUS ACTION: None
\\~\\
\textcolor{red}{example\_assistant} 

\texttt{\`{}\`{}\`{}}stop [\$279.49]\texttt{\`{}\`{}\`{}}
\\~\\

\textcolor{red}{example\_user}

OBSERVATION:

[164] textbox 'Search' focused: True required: False

[171] button 'Go'

[174] link 'Find directions between two points'

[212] heading 'Search Results'

[216] button 'Close'

URL: http://openstreetmap.org

OBJECTIVE: Show me the restaurants near ABC

PREVIOUS ACTION: None
\\~\\
\textcolor{red}{example\_assistant}

\texttt{\`{}\`{}\`{}}type [164] [restaurants near ABC] [1]\texttt{\`{}\`{}\`{}}
        }%
}
\caption{The two examples provided as \texttt{example\_user} and \texttt{example\_assistant} for the direct agent. The agent directly emits the next action given the observation.}
    \label{fig:direct_example}
\end{figure}

\subsection{Additional Error Analysis}\label{sec:additional_error_analysis}
\begin{figure}
    \centering
    \includegraphics[width=\linewidth]{figs/errors_1.pdf}
    \caption{Two examples where the \textsc{GPT-4} agent failed, along with their screenshot and the accessibility tree of the relevant sections (grey). On the left, the agent fails to proceed to the ``Users'' section to accomplish the task of ``Fork all Facebook repos''; on the right, the agent repeats entering the same search query even though the observation indicates the input box is filled.}
    \label{fig:additional_error}
\end{figure}

\paragraph{Observation Bias} Realistic websites frequently present information on similar topics across various sections to ensure optimal user accessibility. However, a \textsc{GPT-4} agent often demonstrates a tendency to latch onto the first related piece of information it encounters without sufficiently verifying its relevance or accuracy. 
For instance, the homepage of the E-Commerce CMS displays the best-selling items based on \emph{recent purchases}, while historical best-seller data is typically accessed via a separate report. Presented with the task of \sent{What is the top-1 best-selling product in 2022}, the \textsc{GPT-4} agent defaults to leveraging the readily available information on the homepage, bypassing the necessary step of generating the report to obtain the accurate data.

\paragraph{Failures in Observation Interpretation} Interestingly, while \textsc{GPT-4} is capable of summarizing the observations, it occasionally overlooks more granular information, such as the previously entered input. As in the right-hand example of \autoref{fig:additional_error}, \texttt{[5172] StaticText} indicates that the search term ``DMV area'' has already been entered. However, the agent disregards this detail and continuously issues the command \texttt{type [2430] [DMV area]} until it reaches the maximum step limit. Furthermore, the agent often neglects the previous action information that is provided alongside the observation.

We hypothesize that these observed failures are related to the current pretraining and supervised fine-tuning on dialogues employed in GPT models~\cite{ouyang2022training}. These models are primarily trained to execute instructions given \emph{immediate} observations (\ie, the dialogue history); thereby, they may exhibit a lack of explorations. 
Furthermore, in dialogue scenarios, subtle differences in NL expressions often have less impact on the overall conversation. As a result, models may tend to overlook minor variations in their observations.

% \subsection{A Successful Trajectory}\label{app:succ_example}
% A successful trajectory is shown in \autoref{fig:success_example}.
% \begin{figure*}
%     \centerline{\includegraphics[width=1.25\linewidth]{figs/s1.pdf}}
% \centerline{\includegraphics[width=1.25\linewidth]{figs/s2.pdf}}
% \end{figure*}

% \begin{figure*}
% \centerline{\includegraphics[width=1.25\linewidth]{figs/s3.pdf}}
% \centerline{\includegraphics[width=1.25\linewidth]{figs/s4.pdf}}
%     \caption{The trajectory of the \textsc{GPT-4} reasoning agent on performing the task of ``Tell me the status of my latest order and when will it arrive''. The agent uses accessibility tree as the web page observations, however, for readibility, we display the procedure using screenshots. In each screenshot, the element predicted for interaction by the agent is highlighted by a red arrow. 
%     The agent's reasoning process and predicted actions are beneath each screenshot.}
%     \label{fig:success_example}
% \end{figure*}